\section{Π/Σ-types and the calculus of constructions}\label{sec:01-basics:05-pi-sigma-and-coc:1}

§3 built Π-types concretely, from \texttt{Fin n} and \texttt{Vec.replicate}, up to the general pattern \(\prod_{x:A} B(x)\). This section adds the second half of the picture --- \textbf{Σ-types}, the dependent pair dual to Π --- and then zooms out to name the whole formal system these pieces belong to: the \textbf{calculus of constructions} (CoC), the core type theory underlying Lean. (Universes get their \emph{own} formal typing rule in \hyperref[sec:05-rigor-check:03-typing-rules-and-safety:1]{Chapter 5 §3}, once Chapter 5 has built up the informal picture of \texttt{Type}/\texttt{Type 1} those rules formalize --- nothing here depends on having read that yet.)

\subsection{Π-types, briefly recalled}\label{sec:01-basics:05-pi-sigma-and-coc:2}

\[
\prod_{x : A} B(x)
\]

--- ``a function that, given \(x : A\), returns a term of type \(B(x)\), a type allowed to mention \(x\).'' §3's \texttt{Vec.replicate}, with type \texttt{(n : Nat) → Vec α n}, literally \emph{is} \(\prod_{n:\mathtt{Nat}} \mathrm{Vec}\,\alpha\,n\); Lean's surface syntax \texttt{(x : A) → B x} \textbf{is} \(\prod_{x:A} B(x)\), and \texttt{∀} is the same construct specialized to \texttt{B : A → Prop}. Every \texttt{∀ n : Nat, P n} since Chapter 3 was already a Π-type.

\begin{progcorner}

Even Python's \texttt{TypeVar} (its own fix for generic functions like \texttt{identity}) cannot express \texttt{Vec.replicate}'s type:

\end{progcorner}

\begin{lstlisting}[language=Python,style=python]
from typing import TypeVar
T = TypeVar("T")

def identity(x: T) -> T:      # generic over a TYPE — fine, this is TypeVar's job
    return x

# There is no Python type-hint equivalent of:
#   def replicate(a: T, n: <this specific Nat>) -> Vec[T, <that same Nat>]
# because "the type mentions n's VALUE" is not expressible by any
# combination of ordinary hints or TypeVar.
\end{lstlisting}

This is the honest answer to why Lean needs a whole extra concept here: it is not that Python's type hints are missing a minor convenience, it is that \emph{no} mainstream static type system --- Python's, Java's, C\#'s, TypeScript's --- has a construct for ``the type mentions a specific runtime value,'' because none of them needed a proof assistant's level of precision. Π-types are precisely that missing construct, made rigorous.

\subsection{\texorpdfstring{\texttt{Prop} as a special, proof-irrelevant universe}{Prop as a special, proof-irrelevant universe}}\label{sec:01-basics:05-pi-sigma-and-coc:3}

Lean's \texttt{Prop} (Chapter 3) sits alongside \texttt{Type} as its own universe, with one extra rule: \textbf{proof irrelevance}. Any two terms of the same type \texttt{P : Prop} are considered definitionally equal, since a proof carries no computational content beyond the bare fact that \emph{a} proof exists:

\begin{lstlisting}
theorem two_proofs (h1 h2 : 2 + 2 = 4) : h1 = h2 := rfl
\end{lstlisting}

\texttt{rfl} succeeds no matter \emph{how} \texttt{h1} and \texttt{h2} were each proved --- by \texttt{rfl} directly, by a long tactic block, by an entirely different chain of lemmas --- because Lean does not distinguish between different proofs of the same \texttt{Prop}. This is what makes Curry--Howard's slogan ``propositions are types, proofs are terms'' more than a slogan: \texttt{P : Prop} really is a type, \texttt{h : P} really is an ordinary term of that type built by ordinary \(\lambda\)/application/Π machinery, and the \emph{only} difference from an ordinary \texttt{Type} is that Lean does not care \emph{which} term of \texttt{P} was produced --- only that one was produced.

\subsection{Σ-types: the dependent pair, dual to Π}\label{sec:01-basics:05-pi-sigma-and-coc:4}

Where Π generalizes \(\to\), a \textbf{Σ-type} (dependent pair / dependent sum) generalizes \(\times\) (the ordinary product):

\[
\sum_{x:A} B(x)
\]

a pair \(\langle a, b\rangle\) with \(a : A\) and \(b : B(a)\) --- the \emph{second} component's type is allowed to depend on the \emph{first} component's \emph{value}.

\textbf{Why ``sum,'' if it generalizes ×?} The name is not a mismatch between Σ and the product --- it names the more fundamental view underneath both. \(\sum_{x:A} B(x)\) is, categorically, a disjoint union (coproduct) of the family \(B\) \emph{indexed by} \(A\): one tagged copy of \(B(x)\) for every \(x \in A\), glued together as separate parts that never overlap. \hyperref[sec:03-propositions-and-proofs:01-prop:1]{Chapter 3 §1}'s Curry--Howard table already lists \texttt{∧} as a \textbf{product type} and \texttt{∨} as a \textbf{sum (coproduct) type}, side by side with \texttt{∃} as a \textbf{Σ-type} --- as though all three were unrelated rows. They are not: \texttt{∧} and \texttt{∨} are both shadows of this one Σ-type construction, cast along two different axes of ``what is allowed to vary'':

\begin{itemize}
\tightlist
\item
  \textbf{Constant family, varying index} --- if \(B(x)\) is the \emph{same} type \(B\) for every \(x\), the disjoint union of \(|A|\) tagged copies of \(B\) coincides with the ordinary product \(A \times B\) (a value of \(A\), tagging \emph{which} copy, paired with a value of \(B\)). This is the ``dependent pair'' reading used throughout this section, and it is exactly Chapter 3's \texttt{∧} (\texttt{P ∧ Q} is \texttt{Σ}-like with the index type restricted to two unlabeled slots, both of type \texttt{Prop}).
\item
  \textbf{Two-point index, varying family} --- if \(A\) is a two-element type (Lean's \texttt{Bool}, or ``which side'' of an \texttt{Or}) and \(B(\mathrm{true}) = P\), \(B(\mathrm{false}) = Q\) for two unrelated propositions, the \emph{same} construction \(\sum_{x:\mathrm{Bool}} B(x)\) collapses instead to \(P \sqcup Q\) --- the ordinary sum/coproduct type, exactly Chapter 3's \texttt{P ∨ Q}, built from \texttt{Or.inl}/\texttt{Or.inr}.
\end{itemize}

Σ is called a \emph{sum} because it always literally is one --- an indexed disjoint union --- and the product reading (\texttt{×}, and \texttt{∧} as its \texttt{Prop} special case) is the constant-family special case, not the general construction. \texttt{∀}/Π sits on the dual side of exactly the same coin: \(\prod_{x:A} B(x)\) is an indexed \emph{product}, which for a constant family collapses to the ordinary function type \(A \to B\) (an exponential, \(B^A\)) rather than to \(A \times B\) --- which is also why Π is never confused with \texttt{×} the way Σ is with \texttt{∨}, and no equivalent puzzle arises on that side.

This shape has already appeared, without the name: §3's \texttt{Fin n} is, under the hood, exactly this pair ---

\begin{lstlisting}
#print Fin
-- structure Fin (n : Nat) : Type
-- fields:
--   Fin.val  : Nat
--   Fin.isLt : val < n
\end{lstlisting}

a \texttt{Nat} value \texttt{val}, paired with a proof \texttt{isLt} whose \emph{statement} (\texttt{val < n}) mentions \texttt{val} itself. That is \(\Sigma\), concretely: the second field's type depends on the value of the first. Here is the general construct, spelled out with Lean's actual \texttt{Sigma}, pairing a bound with an actual number below it:

\begin{lstlisting}
def mySigma : Σ n : Nat, Fin n := ⟨3, ⟨2, by decide⟩⟩
#eval mySigma.fst        -- 3
#eval mySigma.snd.val    -- 2
\end{lstlisting}

\texttt{mySigma.fst} extracts the first component with no fuss at all --- it is ordinary data, just like extracting \texttt{.1} from an ordinary pair. This extractability matters, because it is exactly what \texttt{∃} (next) does \emph{not} allow, and the reason why is the most subtle point in this section.

This dependent-pair shape is also the ``structure bundling data + proofs'' pattern that recurs throughout this book: \texttt{Group G}'s \texttt{⟨op, id, inv, assoc, ...⟩} is, underneath Lean's \texttt{structure} sugar, an iterated Σ-type --- a witness \texttt{op}, paired with a witness \texttt{id}, paired with \texttt{assoc}, whose \emph{type} depends on the values of \texttt{op} and \texttt{id} supplied earlier in the same structure. Chapter 2's ``structures can bundle proofs alongside data'' was already, silently, an appeal to Σ.

\textbf{A caveat about \texttt{∃}, worth being precise about.} Chapter 3 reads \texttt{∃ x : α, P x} as ``a structure: a witness value plus a proof that the witness satisfies \texttt{P}'' --- in effect a Σ-type, with \texttt{P : α → Prop} playing the role of \texttt{B}. Lean's actual \texttt{Exists} is, however, \emph{not literally} the Σ-type \texttt{Sigma}. \texttt{Exists} is specifically built to land in \texttt{Prop}, and by proof irrelevance (above), that means the witness cannot be \emph{extracted} from an \texttt{∃}-proof computationally:

\begin{lstlisting}
def bad (h : ∃ n : Nat, n > 0) : Nat := h.1
\end{lstlisting}

\begin{lstlisting}
error(lean.projNonPropFromProp): Invalid projection: Cannot project a
value of non-propositional type
  Nat
from the expression
  h
which has propositional type
  ∃ n, n > 0
\end{lstlisting}

Compare this directly to \texttt{mySigma.fst} two paragraphs up, which worked with no error at all --- same pair-shape, different universe, and that difference is exactly why one extracts and the other does not. If \texttt{∃} \emph{did} allow extracting a witness, two different (but both equally valid) choices of witness could produce different results from a ``proof-irrelevant'' input, contradicting proof irrelevance itself. \texttt{Sigma} (the actual, \texttt{Type}-valued dependent pair, matching a \texttt{structure}'s bundled data) \emph{does} support projecting out its first component, precisely because it does not carry \texttt{Exists}'s irrelevance guarantee. So: \texttt{∃} has the same shape as Σ but lives in a universe where the witness is unobservable. That is what makes it a \emph{restricted} cousin of Σ rather than literally Σ. The \texttt{structure}-based bundling this book uses throughout for \texttt{Group}/\texttt{Ring}/\texttt{Module} genuinely is \texttt{Sigma}-like (extractable), while an \texttt{∃}-statement genuinely is not.

\subsection{The calculus of constructions, named}\label{sec:01-basics:05-pi-sigma-and-coc:5}

Putting the pieces together, the \textbf{calculus of constructions} (CoC, Coquand--Huet) --- the core type theory underlying Lean, Coq, and similar systems --- is a small formal system built from a variable, a function abstraction, and function application, plus:

\begin{enumerate}
\def\labelenumi{\arabic{enumi}.}
\tightlist
\item
  an infinite hierarchy of type universes \(\mathtt{Type}\,0,
  \mathtt{Type}\,1, \ldots\), formalized in \hyperref[sec:05-rigor-check:03-typing-rules-and-safety:1]{Chapter 5 §3},
\item
  Π-types generalizing \(\to\), allowing types to depend on terms (this chapter's §3 and the top of this page),
\item
  (in Lean's specific extension, the \textbf{Calculus of Inductive Constructions}, CIC) \texttt{inductive} type declarations --- \texttt{Nat}, \texttt{Bool}, \texttt{Vec}/\texttt{Fin} (§3), \texttt{Path} (Chapter 11), and every \texttt{structure} written throughout --- giving Σ-types and much more (arbitrary recursive data) beyond what bare CoC provides,
\item
  \texttt{Prop} as a distinguished, proof-irrelevant universe (Chapter 3, and above), making Curry--Howard a precise correspondence rather than an analogy.
\end{enumerate}

Every single Lean construct used in this book --- \texttt{def}, \texttt{structure}, \texttt{theorem}, \texttt{∀}, \texttt{∃}, tactics (which merely \emph{build} CIC terms, one piece at a time, without the terms being written by hand) --- compiles down to a term in exactly this system, checked by Lean's kernel using nothing more than typing rules like the ones sketched here and in Chapter 5, applied mechanically. \hyperref[sec:05-rigor-check:03-typing-rules-and-safety:1]{Chapter 5 §3} completes the picture with the simply typed λ-calculus's typing judgments and its progress/preservation theorems --- the formal reason ``well-typed proofs do not go wrong'' --- plus the universe-formation rule only sketched by name above.

\subsection{References}\label{sec:01-basics:05-pi-sigma-and-coc:6}

Full citations in the \hyperref[sec:root:bibliography:1]{Bibliography}.

\begin{itemize}
\tightlist
\item
  Coquand and Huet (\cite{CoquandHuet1988}) --- the original paper defining CoC, the system this section formalizes.
\item
  The Univalent Foundations Program (\cite{HoTT2013}) --- Ch. 1 gives a careful, from-scratch treatment of Π-types, Σ-types, and universes, in notation closely matching this section's.
\item
  Martin-Löf (\cite{MartinLof1984}) --- the foundational source for dependent Π/Σ types and universes predating CoC, for readers wanting the idea in its original, non-CoC-specific form.
\item
  \emph{Theorem Proving in Lean 4} (\cite{TPIL4}), ``Dependent Types'' and ``Propositions and Proofs'' --- Lean's own documentation on \texttt{Prop}, proof irrelevance, and universes, matching the presentation here.
\item
  All Lean code in this section was checked directly against the toolchain pinned in this repository's \texttt{lean\_\allowbreak{}project/\allowbreak{}lean-\allowbreak{}toolchain} rather than only described; the \texttt{\#print Fin} output and both error messages shown are copied from real \texttt{lake env lean} runs.
\end{itemize}
